\Needspace{11\baselineskip}
\SourceLeaf{281}
\begin{center}
\textit{Vergleich der Klassenzahlen in den einzelnen Geschlechtern verschiedener Ordnungen.}
\end{center}

\Article{253}
Der vorige Satz setzt dieselbe Ordnung voraus und darf nicht auf verschiedene Ordnungen ausgedehnt werden. So gibt es \textit{zum Beispiel}\ bei der Determinante $-171$ $20$ positive Klassen, die auf vier Ordnungen verteilt sind: Die eigentlich primitive Ordnung enthält zwei Geschlechter mit jeweils sechs Klassen; in der uneigentlich primitiven Ordnung besitzen zwei Geschlechter vier Klassen, jedes zwei; in der von der eigentlich primitiven Ordnung der Determinante $-19$ abgeleiteten Ordnung gibt es ein einziges Geschlecht mit drei Klassen; schließlich hat die von der uneigentlich primitiven Ordnung der Determinante $-19$ abgeleitete Ordnung ein einziges Geschlecht, das aus einer Klasse besteht. Ebenso verhält es sich mit den negativen Klassen. Es lohnt sich also, nach einem allgemeinen Prinzip zu suchen, von dem der Zusammenhang zwischen den Klassenzahlen verschiedener Ordnungen abhängt. Wir nehmen an, $K,L$ seien zwei Klassen derselben (positiven) Ordnung $O$ der Determinante $D$ und $M$ eine eigentlich primitive Klasse derselben Determinante, aus deren Komposition mit $K$ die Klasse $L$ entsteht; eine solche lässt sich nach Artikel~251 stets angeben. In manchen Fällen kann $M$ die \textit{einzige} eigentlich primitive Klasse sein, die mit $K$ zusammengesetzt $L$ ergibt; in anderen können mehrere verschiedene eigentlich primitive Klassen diese Eigenschaft haben. Wir nehmen allgemein an, es gebe $r$ solche eigentlich primitiven Klassen $M,M',M''\ldots M^{r-1}$, die jeweils mit $K$ zusammengesetzt dieselbe Klasse $L$ ergeben, und bezeichnen ihre Menge mit $W$. Ferner sei $L'$ eine andere Klasse der Ordnung $O$ (verschieden von der Klasse $L$) und $N'$ eine eigentlich primitive Klasse der Determinante $D$, die mit $L$ zusammengesetzt $L'$ ergibt. Die Menge der Klassen $N'+M,N'+M',N'+M''\ldots N'+M^{r-1}$, die sämtlich eigentlich primitiv und voneinander verschieden sind, werde mit $W'$ bezeichnet. Dann erkennt man leicht, dass $K$ mit irgendeiner Klasse aus $W'$ zusammengesetzt $L'$ ergibt; daraus folgt, dass $W$ und $W'$ keine Klasse gemeinsam haben. Außerdem bestätigt man ohne Mühe, dass es keine eigentlich primitive Klasse außerhalb der Menge $W'$ gibt, die mit $K$ zusammengesetzt $L'$ ergibt. Ebenso ist klar: Ist $L''$ eine weitere Klasse der Ordnung $O$, verschieden von den Klassen $L,L'$, so gibt es $r$ eigentlich primitive Formen, die sowohl untereinander als auch von den Formen $W,W'$ verschieden sind und jeweils mit $K$ zusammengesetzt $L''$ ergeben; ebenso verhält es sich bei allen übrigen Klassen der Ordnung $O$. Da aber jede eigentlich primitive (positive) Klasse der Determinante $D$ mit $K$ zusammengesetzt eine Klasse der Ordnung $O$ ergibt, folgt leicht: Ist die Gesamtzahl der Klassen der Ordnung $O$ gleich $n$, so ist die Gesamtzahl aller eigentlich primitiven (positiven) Klassen derselben Determinante gleich $rn$.

\SourceLeaf{282}Wir erhalten also die allgemeine Regel: Bezeichnen \(K\), \(L\) beliebige Klassen der Ordnung \(O\) und \(r\) die Anzahl verschiedener eigentlich primitiver Klassen derselben Determinante, die jeweils mit \(K\) zusammengesetzt \(L\) ergeben, so ist die Gesamtzahl aller Klassen in der eigentlich primitiven (positiven) Ordnung \(r\)-mal so groß wie die Anzahl der Klassen der Ordnung \(O\).

Da die Klassen \(K\), \(L\) in der Ordnung \(O\) völlig beliebig gewählt werden können, dürfen wir auch identische Klassen nehmen. Dabei ist es zweckmäßig, diejenige Klasse zu verwenden. in der die einfachste Form dieser Ordnung liegt. Nehmen wir sie also für \(K\) und \(L\), so ist die Aufgabe darauf zurückgeführt, alle eigentlich primitiven Klassen anzugeben, die mit \(K\) zusammengesetzt \(K\) selbst wieder ergeben. Den Weg dazu bereitet der folgende Satz.

\Needspace{8\baselineskip}
\Article{254}
\textsc{Satz.} \textit{Ist \(F=(A,B,C)\) die einfachste Form der Ordnung \(O\) der Determinante \(D\) und \(f=(a,b,c)\) eine eigentlich primitive Form derselben Determinante, so lässt sich durch diese Form \(f\) die Zahl \(AA\) darstellen, wenn \(F\) durch Komposition der Formen \(f\), \(F\) entsteht; und umgekehrt ist \(F\) aus sich selbst und \(f\) zusammengesetzt, wenn \(AA\) durch \(f\) dargestellt werden kann.}

\textit{Beweis.} I. Geht \(F\) in das Produkt \(fF\) durch die Substitution \(p,\allowbreak p',\allowbreak p'',\allowbreak p''';\allowbreak \ q,\allowbreak q',\allowbreak q'',\allowbreak q'''\) über, so ergibt Artikel 235
\[
\begin{gathered}A(aq''q''-2bqq''+cqq)=A^3,\quad\text{also}\\ AA=aq''q''-2bqq''+cqq.\quad\textit{Erster Teil bewiesen.}\end{gathered}
\]

II. Es werde \textit{angenommen}, dass \(AA\) durch \(f\) darstellbar ist. Die Werte der Unbestimmten, durch die dies geschieht, bezeichnen wir mit \(q''\). \(-q\), also sei \(AA=aq''q''-2bqq''+cqq\), und setzen
\[
\begin{aligned}
q''a-q(b+B)&=Ap,\\ -qC&=Ap',\\ q''(b-B)-qc&=Ap''\\ -q''C&=Ap''',\\ q''a-q(b-B)&=Aq',\\ q''(b+B)-qc&=Aq'''
\end{aligned}
\]
Dann bestätigt man leicht, dass \(F\) in das Produkt \(fF\) durch die Substitution \(p,\allowbreak p',\allowbreak p'',\allowbreak p''';\allowbreak \ q,\allowbreak q',\allowbreak q'',\allowbreak q'''\) übergeht. und somit aus \(f\) und \(F\) zusammengesetzt ist, sofern alle Zahlen \(p\), \(p'\) usw. ganzzahlig sind. Nach der Beschreibung der einfachsten Form ist nun \(B\) entweder \(0\) oder \(\frac12 A\), also \(\frac{2B}{A}\) ganzzahlig; daraus ist auch klar, dass \(\frac{C}{A}\) stets ganzzahlig ist. Somit sind \(q'-p\), \(p'\), \(q'''-p''\), \(p'''\) ganzzahlig, und es bleibt nur noch zu beweisen, dass \(p\) und \(p''\) ganzzahlig sind. Es gilt aber
\SourceLeaf{283}
\[
pp+\frac{2pqB}{A}=a-\frac{qqC}{A},\quad p''p''+\frac{2p''q''B}{A}=c-\frac{q''q''C}{A}
\]
also gilt im Fall \(B=0\)
\[
pp=a-\frac{qqC}{A},\quad p''p''=c-\frac{q''q''C}{A}
\]
und daher sind \(p,p''\) ganzzahlig; gilt dagegen \(B=\frac12 A\), so ist
\[
pp+pq=a-\frac{qqC}{A},\quad p''p''+p''q''=c-\frac{q''q''C}{A}
\]
woraus ebenso leicht folgt, dass \(p\) und \(p''\) auch in diesem Fall ganzzahlig sind. Daraus schließen wir, dass \(F\) aus \(f\) und \(F\) zusammengesetzt ist. \textit{Zweiter Teil bewiesen.}

\Needspace{8\baselineskip}
\Article{255}
Die Aufgabe ist also darauf zurückgeführt, alle eigentlich primitiven Klassen der Determinante \(D\) anzugeben, durch deren Formen \(AA\) darstellbar ist. Offenbar kann \(AA\) durch jede Form dargestellt werden, deren erster Koeffizient entweder \(AA\) oder das Quadrat eines Teilers von \(A\) ist. Kann umgekehrt \(AA\) durch eine Form \(f\) dargestellt werden, indem man ihren Unbestimmten die Werte \(\alpha e\), \(\gamma e\) mit dem größten gemeinsamen Teiler \(e\) gibt, so geht die Form \(f\) durch die Substitution \(\alpha,\beta,\gamma,\delta\) in eine Form über, deren erster Koeffizient \(\frac{AA}{ee}\) ist; diese Form ist der Form \(f\) eigentlich äquivalent, wenn \(\beta\), \(\delta\) so gewählt werden, dass \(\alpha\delta-\beta\gamma=1\). Daher gibt es in jeder Klasse, durch deren Formen \(AA\) darstellbar ist, Formen, deren erster Koeffizient \(AA\) oder das Quadrat eines Teilers von \(A\) ist. Es kommt also darauf an, alle eigentlich primitiven Klassen der Determinante \(D\) zu finden, in denen solche Formen vorkommen. Das geschieht folgendermaßen: Es seien \(a,a',a''\) usw. alle (positiven) Teiler von \(A\); man suche alle Werte des Ausdrucks \(\surd D\ (\mathrm{mod.}\,aa)\), die einschließlich der Grenzen zwischen \(0\) und \(aa-1\) liegen, bezeichne sie mit \(b,b',b''\) usw. und setze
\[
bb-D=aac,\quad b'b'-D=aac',\quad b''b''-D=aac''\text{ etc.}
\]
Die Menge der Formen \((aa,b,c)\), \((aa,b',c')\) usw. werde mit \(V\) bezeichnet. Dann erkennt man leicht, dass jede Klasse der Determinante \(D\), in der eine Form mit dem ersten Koeffizienten \(aa\) vorkommt, auch eine Form aus \(V\) enthalten muss. Ebenso suche man alle Formen der Determinante \(D\), deren erster Koeffizient \(a'a'\) ist und deren mittlerer Koeffizient einschließlich der Grenzen zwischen \(0\) und \(a'a'-1\) liegt, und bezeichne ihre Menge mit \(V'\); entsprechend sei \(V''\)
\SourceLeaf{284}
die Menge solcher Formen mit dem ersten Koeffizienten \(a''a''\) usw. Aus \(V,V',V''\) usw. entferne man alle Formen, die nicht eigentlich primitiv sind, ordne die übrigen in Klassen ein und behalte, falls mehrere zu derselben Klasse gehören, in jeder Klasse nur eine. So erhält man alle gesuchten Klassen; ihre Anzahl verhält sich zur Einheit wie die Anzahl aller eigentlich primitiven (positiven) Klassen zur Anzahl der Klassen in der Ordnung \(O\).

\textit{Beispiel.} Es sei \(D=-531\) und \(O\) die positive Ordnung, die von der uneigentlich primitiven Ordnung der Determinante \(-59\) abgeleitet ist; ihre einfachste Form ist \((6,3,90)\), also \(A=6\). Hier sind \(a,a',a'',a'''\) gleich \(1,2,3,6\); \(V\) enthält die Form \((1,0,531)\); \(V'\) die Formen \((4,1,133)\), \((4,3,135)\); \(V''\) die Formen \((9,0,59)\), \((9,3,60)\), \((9,6,63)\); schließlich \(V'''\) die Formen \((36,3,15)\), \((36,9,17)\), \((36,15,21)\), \((36,21,27)\), \((36,27,35)\), \((36,33,45)\). Von diesen zwölf Formen müssen aber sechs entfernt werden, nämlich aus \(V''\) die zweite und dritte, aus \(V'''\) die erste, dritte, vierte und sechste; diese sind sämtlich abgeleitete Formen. Die sechs übrigen gehören, wie man findet, alle zu verschiedenen Klassen. Tatsächlich ist die Anzahl der eigentlich primitiven (positiven) Klassen der Determinante \(-531\) gleich \(18\), und die Anzahl der uneigentlich primitiven (positiven) Klassen der Determinante \(-59\) (oder die Anzahl der Klassen der Determinante \(-531\), die von diesen abgeleitet sind) gleich \(3\); somit verhält sich die erste Anzahl zur zweiten wie \(6\) zu \(1\).
